\documentclass[a4paper,10pt]{report}\input{../0.Préambule/Préambule_cours_prof}\begin{document}

\pagestyle{empty}
\newpage
\noindent NOM, Prénom et classe : \dotfill

\section*{Expressions littérales - Conversion \hspace{3cm} Sujet A}
\addcontentsline{toc}{section}{Expressions littérales - Prismes droit \hspace{3cm} Sujet A}

\noindent \textit{L'usage de la calculatrice est \textbf{autorisée}.} \hfill \textit{Durée : 55 minutes}

\paragraph{Exercice 1} : \hfill \textbf{2 points} 

\noindent Donner les différentes situations dans lesquelles il est possible de ne pas écrire le symbole \og $\times$ \fg{}

\paragraph{Exercice 2} : \hfill \textbf{2,5 points}

\noindent Simplifie les expressions suivantes :

\vspace{-2mm}
\begin{enumerate}
\begin{tabularx}{\linewidth}{*{5}{>{\arraybackslash}X}}
\item $2 \times a + 5 \times c$ &
%\item $a \times d + 5 \times 8$ &
\item $38 \times (3 + 2 \times c)$ &
%\item $3 \times z - 0 \times b$ &
\item $3 \times 7 - d \times b$ &
%\item $a \times (3 \times 9 + b \times n)$ &
\item $0 \times u + 1 \times m$ &
%\item $a \times 6 \times n + 3 \times p$ &
\item $a \times n \times n + p \times p \times p$ \\
%\item $9 \times m \times 5 + k \times j \times 8$ \\
\end{tabularx}
\end{enumerate}

\paragraph{Exercice 3} : \textit{A faire sur la fiche} \hfill \textbf{3,5 points}

\noindent Relie chaque phrase de gauche à l'expression littérale correspondante de droite.

\renewcommand{\arraystretch}{1.1}
\begin{tabularx}{14cm}{r c X c l}
Somme de $y$ et de $7$ & $\bullet$ & & $\bullet$ & $7 \times (y-3)$ \\
Produit de $7$ par la somme de $y$ et de $3$ & $\bullet$ & & $\bullet$ & $7-y$ \\
Produit de $7$ par la différence entre $y$ et $3$ & $\bullet$ & & $\bullet$ & $y+7\times 3$ \\
Différence du produit de $7$ par $y$ et de $3$ & $\bullet$ & & $\bullet$ & $y+7$ \\
Différence entre $7$ et $y$ & $\bullet$ & & $\bullet$ & $7 \times y + 3$ \\
Somme de $y$ et du produit de $3$ par $7$ & $\bullet$& & $\bullet$ & $7 \times (y+3)$ \\
Somme du produit de $7$ par $y$ et de $3$ & $\bullet$& & $\bullet$ & $7 \times y - 3$ \\
\end{tabularx}

\paragraph{Exercice 4} : \hfill \textbf{3 points}

Dans une assemblée, il y a deux fois plus de Belges que de Luxembourgeois et $48$ Néerlandais de plus que de Luxembourgeois. 

\noindent On désigne par $x$ le nombre de Luxembourgeois. 

\noindent Quelle est la composition de l'assemblée (en fonction de $x$).

\paragraph{Exercice 5} : \hfill \textbf{3 points}

\vspace{5mm}
\noindent \begin{minipage}{6cm}
\arrayrulecolor{black}
\renewcommand{\arraystretch}{0.1}
\begin{tabular}{|p{5.5cm}|}
\hline
\vspace{2mm} 1. \edupythoninline{Choisis un nombre} \\ 
\vspace{2mm} 2. \edupythoninline{Calcule son double} \\ 
\vspace{2mm} 3. \edupythoninline{Ajoute 8} \\ 
\vspace{2mm} 4. \edupythoninline{Divise le resultat par 2} \\ 
\vspace{2mm} 4. \edupythoninline{Soustrait le nombre de depart} \\ 
\vspace{2mm} \\
\hline
\end{tabular}
\end{minipage}
\hspace{5mm}
\begin{minipage}{10.2cm}
\begin{enumerate}
\item Effectue le programme de calcul pour $5$.
\item Recommence avec deux autres nombres. Que constates-tu ?
\item Soit $x$ le nombre choisi à l'étape $1$. 

Exprime le résultat du programme de calcul en fonction de $x$.
\end{enumerate}
\end{minipage}

\paragraph{Exercice 6} : \hfill \textbf{3 points}

\noindent Calcule la valeur de $M$, de $R$ et de $E$ pour $m=6$ et $n=7$

\begin{enumerate}
\noindent \begin{tabularx}{\linewidth}{*{3}{>{\arraybackslash}X}}
\item $M = 7m + 8n + nm$&
\item $R = 8n + 5mn -8n$&
\item $E = 8n - 4m - 8mn$ \\
\end{tabularx}
\end{enumerate}

\paragraph{Exercice 7} : \textit{A faire sur la fiche} \hfill \textbf{3 points}

\begin{enumerate}
\item Effectue les conversions suivantes.

\begin{enumerate}
\begin{tabularx}{\linewidth}{*{3}{>{\arraybackslash}X}}
%\item $1$dm$^3$ = \hspace{1.5cm} L &
%\item $1$m$^3$ = \hspace{1.5cm} L &
%\item $1$mL = \hspace{1.5cm} cm$^3$ \\
\item $232,4$L = \hspace{1.5cm} m$^3$ &
\item $56,78$cm$^3$ = \hspace{1.5cm} dL &
\item $\np{7302}$L = \hspace{1.5cm} mm$^3$ \\
\end{tabularx}
\end{enumerate}

\item Complète avec l'unité de capacité (L, dL, ...) qui convient.

\begin{enumerate}
\begin{tabularx}{\linewidth}{*{3}{>{\arraybackslash}X}}
\item $\np{26}$dm$^3$ = $26$ &
\item $\np{0,502}$dm$^3$ = $\np{502}$ &
\item $\np{2}$m$^3$ = $\np{2000}$ \\
%\item $\np{7,8}$cm$^3$ = $0,78$ &
%\item $\np{0,17}$dam$^3$ = $\np{17000}$ &
%\item $\np{3542}$mm$^3$ = $\np{3,542}$ \\
\end{tabularx}
\end{enumerate}
\end{enumerate}

\newpage
\noindent NOM, Prénom et classe : \dotfill

\section*{Expressions littérales - Conversion \hspace{3cm} Sujet B}
\addcontentsline{toc}{section}{Expressions littérales - Prismes droit \hspace{3cm} Sujet B}

\noindent \textit{L'usage de la calculatrice est \textbf{autorisée}.} \hfill \textit{Durée : 55 minutes}

\paragraph{Exercice 1} : \hfill \textbf{2 points} 

\noindent Donner les différentes situations dans lesquelles il est possible de ne pas écrire le symbole \og $\times$ \fg{}

\paragraph{Exercice 2} : \hfill \textbf{2,5 points}

\noindent Simplifie les expressions suivantes :

\vspace{-2mm}
\begin{enumerate}
\begin{tabularx}{\linewidth}{*{5}{>{\arraybackslash}X}}
%\item $2 \times a + 5 \times c$ &
\item $a \times d + 5 \times 8$ &
%\item $38 \times (3 + 2 \times c)$ &
\item $3 \times z - 0 \times b$ &
%\item $3 \times 7 - d \times b$ &
\item $a \times (3 \times 9 + b \times n)$ &
%\item $0 \times u + 1 \times m$ &
\item $a \times n + p \times p \times p$ &
%\item $a \times 6 \times n + 3 \times p$ \\
\item $9 \times m \times 5 + k \times j \times 8$ \\
\end{tabularx}
\end{enumerate}

\paragraph{Exercice 3} : \textit{A faire sur la fiche}  \hfill \textbf{3,5 points}

\noindent Relie chaque phrase de gauche à l'expression littérale correspondante de droite.

\renewcommand{\arraystretch}{1.1}
\begin{tabularx}{14cm}{r c X c l}
Différence entre $7$ et $y$ & $\bullet$ & & $\bullet$ & $7 \times y + 3$ \\
Somme de $y$ et du produit de $3$ par $7$ & $\bullet$& & $\bullet$ & $7 \times (y+3)$ \\
Somme du produit de $7$ par $y$ et de $3$ & $\bullet$& & $\bullet$ & $7 \times y - 3$ \\
Somme de $y$ et de $7$ & $\bullet$ & & $\bullet$ & $7 \times (y-3)$ \\
Produit de $7$ par la somme de $y$ et de $3$ & $\bullet$ & & $\bullet$ & $7-y$ \\
Produit de $7$ par la différence entre $y$ et $3$ & $\bullet$ & & $\bullet$ & $y+7\times 3$ \\
Différence du produit de $7$ par $y$ et de $3$ & $\bullet$ & & $\bullet$ & $y+7$ \\
\end{tabularx}

\paragraph{Exercice 4} : \hfill \textbf{3 points}

Dans une assemblée, il y a trois fois plus de Belges que de Luxembourgeois et $36$ Néerlandais de plus que de Luxembourgeois.

\noindent On désigne par $x$ le nombre de Luxembourgeois. 

\noindent Quelle est la composition de l'assemblée (en fonction de $x$).

\paragraph{Exercice 5} : \hfill \textbf{3 points}

\vspace{5mm}
\noindent \begin{minipage}{6cm}
\arrayrulecolor{black}
\renewcommand{\arraystretch}{0.1}
\begin{tabular}{|p{5.5cm}|}
\hline
\vspace{2mm} 1. \edupythoninline{Choisis un nombre} \\ 
\vspace{2mm} 2. \edupythoninline{Calcule son double} \\ 
\vspace{2mm} 3. \edupythoninline{Ajoute 6} \\ 
\vspace{2mm} 4. \edupythoninline{Divise le resultat par 2} \\ 
\vspace{2mm} 4. \edupythoninline{Soustrait le nombre de depart} \\ 
\vspace{2mm} \\
\hline
\end{tabular}
\end{minipage}
\hspace{5mm}
\begin{minipage}{10.2cm}
\begin{enumerate}
\item Effectue le programme de calcul pour $4$.
\item Recommence avec deux autres nombres. Que constates-tu ?
\item Soit $x$ le nombre choisi à l'étape $1$. 

Exprime le résultat du programme de calcul en fonction de $x$.
\end{enumerate}
\end{minipage}

\paragraph{Exercice 6} : \hfill \textbf{3 points}

\noindent Calcule la valeur de $M$, de $R$ et de $E$ pour $m=7$ et $n=6$

\begin{enumerate}
\noindent \begin{tabularx}{\linewidth}{*{3}{>{\arraybackslash}X}}
\item $M = 7m + 10n + nm$&
\item $R = 10n + 5mn -8n$&
\item $E = 8n - 4m - 6mn$ \\
\end{tabularx}
\end{enumerate}

\paragraph{Exercice 7} : \textit{A faire sur la fiche} \hfill \textbf{3 points}

\begin{enumerate}
\item Effectue les conversions suivantes.

\begin{enumerate}
\begin{tabularx}{\linewidth}{*{3}{>{\arraybackslash}X}}
\item $1$dm$^3$ = \hspace{1.5cm} L &
\item $1$m$^3$ = \hspace{1.5cm} L &
\item $1$mL = \hspace{1.5cm} cm$^3$ \\
%\item $232,4$L = \hspace{1.5cm} m$^3$ &
%\item $56,78$cm$^3$ = \hspace{1.5cm} dL &
%\item $\np{7302}$L = \hspace{1.5cm} mm$^3$ \\
\end{tabularx}
\end{enumerate}

\item Complète avec l'unité de capacité (L, dL, ...) qui convient.

\begin{enumerate}
\begin{tabularx}{\linewidth}{*{3}{>{\arraybackslash}X}}
%\item $\np{26}$dm$^3$ = $26$ &
%\item $\np{0,502}$dm$^3$ = $\np{502}$ &
%\item $\np{2}$m$^3$ = $\np{2000}$ \\
\item $\np{7,8}$cm$^3$ = $0,78$ &
\item $\np{0,17}$m$^3$ = $\np{17000}$ &
\item $\np{3542}$mm$^3$ = $\np{3,542}$ \\
\end{tabularx}
\end{enumerate}
\end{enumerate}

\end{document}
